\section{Conclusions}\label{sec:conclusion}
Light sharing, the daily choice between maximum electricity and letting light through to the crop, is the main lever of tracker-based agrivoltaics under a food-security floor. In a differentiable simulator of trackers, PV output, crop light, microclimate and crop growth, evaluated fairly on 400 crop--site pairs with the array density of each site chosen from its own history, light-sharing rules allowed about three times the module density of static arrays at a mean yield retention of 90\%. They raised the land-equivalent ratio from 1.07 to 1.18--1.24, and a rule that follows crop phenology and needs no weather forecast was the best of them. The simulated shade response is conservative against the field trials we could use, so these gains are lower bounds, and calibrating the response raised them without changing the ranking of methods. A learned controller added a further 0.01--0.05 LER, by following a smooth phenology-dependent schedule that a rule can express; training it on an ensemble of shade responses or giving it a farm-measurable biomass signal added nothing. The part where learning through the differentiable simulator paid off was the calibration of the shade response, which gave an ensemble of plausible responses; designing the array density so that the floor holds in 75\% of them restored the compliance of an oracle design at a cost of 0.02--0.03 LER. The yield floor is met on average but not in every season or region, and the electricity potential of 5\% of staple cropland (about 14\,PWh\,yr$^{-1}$) requires about 10\,TWp of modules. Future work should validate the crop response to shade under trackers with more field data, notably for maize and legumes where the meta-analysis and the trials disagree, extend the crop model to grain-number formation and canopy energy balance, add costs and grid constraints, and pursue controllers that guarantee the food-security floor per site and per season.
